
\documentclass{article}
\usepackage{colortbl}
\usepackage{makecell}
\usepackage{multirow}
\usepackage{supertabular}

\begin{document}

\newcounter{utterance}

\centering \large Interaction Transcript for game `dond', experiment `semi', episode 1 with Qwen3.6{-}35B{-}A3B{-}FP8{-}without{-}reasoning.
\vspace{24pt}

{ \footnotesize  \setcounter{utterance}{1}
\setlength{\tabcolsep}{0pt}
\begin{supertabular}{c@{$\;$}|p{.15\linewidth}@{}p{.15\linewidth}p{.15\linewidth}p{.15\linewidth}p{.15\linewidth}p{.15\linewidth}}
    \# & $\;$A & \multicolumn{4}{c}{Game Master} & $\;\:$B\\
    \hline

    \theutterance \stepcounter{utterance}  
    & & \multicolumn{4}{p{0.6\linewidth}}{
        \cellcolor[rgb]{0.9,0.9,0.9}{
            \makecell[{{p{\linewidth}}}]{
                \texttt{\tiny{[P1$\langle$GM]}}
                \texttt{Grasz w grę negocjacji kooperacyjnej, w której musisz uzgodnić z innym graczem sposób podziału zestawu przedmiotów między sobą.} \\
\\ 
\texttt{Zasady:} \\
\texttt{(a) Ty i drugi gracz otrzymujecie zestaw przedmiotów. Każdemu z Was przydzielana jest również tajna funkcja wartości, reprezentująca, ile warte są dla Ciebie poszczególne typy przedmiotów.} \\
\texttt{(b) Wymieniasz wiadomości z drugim graczem, aby uzgodnić, kto otrzyma jakie przedmioty. Każdy może wysłać maksymalnie 5 wiadomości lub skończyć grę wcześniej, składając tajną propozycję w dowolnym momencie.} \\
\texttt{(c) Poproszony o złożenie tajnej propozycji wskazującej przedmioty, które chcesz, sformatowanej w nawiasach kwadratowych w następujący sposób: "{[}Propozycja: <liczba> <nazwa przedmiotu>, <liczba> <nazwa przedmiotu>, <...>{]}"} \\
\texttt{(d) Jeśli Twoje propozycje są się się nawzajem uzupełniają, tzn. jest wystarczająco przedmiotów, aby spełnić obie propozycje, każdy gracz otrzymuje wynik na podstawie sumy wartości przedmiotów, które otrzymał. W przeciwnym razie obaj otrzymujecie zero punktów.} \\
\texttt{(e) Twój cel to maksymalizacja sumy Twojego wyniku i wyniku drugiego gracza.} \\
\\ 
\texttt{Zacznijmy.} \\
\\ 
\texttt{Zestaw dostępnych przedmiotów to:} \\
\\ 
\texttt{2 długopisy, 2 cytryno, 2 garnki} \\
\\ 
\texttt{Twoja tajna funkcja wartości to:} \\
\\ 
\texttt{długopis: 2, cytryna: 0, garnek: 3} \\
\\ 
\texttt{WAŻNE: Twoje wiadomości, chyba że jest to tajna propozycja, są przesyłane bezpośrednio drugiemu graczowi, więc nie dodawaj żadnej odpowiedzi na zasady ani tekstu ogłaszającego Twoją wiadomość. Aby złożyć tajną propozycję, użyj wskazanego formatu. Nie używaj nawiasów kwadratowych podczas komunikowania się z drugim graczem, w przeciwnym razie zostanie to zinterpretowane jako Twoja tajna propozycja.} \\
            }
        }
    }
    & \\ \\

    \theutterance \stepcounter{utterance}  
    & \multicolumn{4}{p{0.6\linewidth}}{
        \cellcolor[rgb]{0.9,0.9,0.9}{
            \makecell[{{p{\linewidth}}}]{
                \texttt{\tiny{[P1$\rangle$GM]}}
                \texttt{Cześć! Widzę, że mamy do podziału 2 długopisy, 2 cytryny i 2 garnki.} \\
\\ 
\texttt{Zanim przejdziemy do konkretów, chciałbym się dowiedzieć, jakie przedmioty są dla Ciebie najbardziej wartościowe? To pomoże nam znaleźć rozwiązanie, które będzie korzystne dla obu nas.} \\
\\ 
\texttt{Moja wiadomość: Jakie są Twoje priorytety w podziale przedmiotów?} \\
            }
        }
    }
    & & \\ \\

    \theutterance \stepcounter{utterance}  
    & & & \multicolumn{2}{p{0.3\linewidth}}{
        \cellcolor[rgb]{0.9,0.9,0.9}{
            \makecell[{{p{\linewidth}}}]{
                \texttt{\tiny{[GM$|$GM]}}
                \texttt{continue} \\
            }
        }
    }
    & & \\ \\

    \theutterance \stepcounter{utterance}  
    & & & \multicolumn{2}{p{0.3\linewidth}}{
        \cellcolor[rgb]{0.9,0.9,0.9}{
            \makecell[{{p{\linewidth}}}]{
                \texttt{\tiny{[GM$|$GM]}}
                \texttt{Cześć! Widzę, że mamy do podziału 2 długopisy, 2 cytryny i 2 garnki.} \\
\\ 
\texttt{Zanim przejdziemy do konkretów, chciałbym się dowiedzieć, jakie przedmioty są dla Ciebie najbardziej wartościowe? To pomoże nam znaleźć rozwiązanie, które będzie korzystne dla obu nas.} \\
\\ 
\texttt{Moja wiadomość: Jakie są Twoje priorytety w podziale przedmiotów?} \\
            }
        }
    }
    & & \\ \\

    \theutterance \stepcounter{utterance}  
    & & \multicolumn{4}{p{0.6\linewidth}}{
        \cellcolor[rgb]{0.9,0.9,0.9}{
            \makecell[{{p{\linewidth}}}]{
                \texttt{\tiny{[GM$\rangle$P2]}}
                \texttt{Grasz w grę negocjacji kooperacyjnej, w której musisz uzgodnić z innym graczem sposób podziału zestawu przedmiotów między sobą.} \\
\\ 
\texttt{Zasady:} \\
\texttt{(a) Ty i drugi gracz otrzymujecie zestaw przedmiotów. Każdemu z Was przydzielana jest również tajna funkcja wartości, reprezentująca, ile warte są dla Ciebie poszczególne typy przedmiotów.} \\
\texttt{(b) Wymieniasz wiadomości z drugim graczem, aby uzgodnić, kto otrzyma jakie przedmioty. Każdy może wysłać maksymalnie 5 wiadomości lub skończyć grę wcześniej, składając tajną propozycję w dowolnym momencie.} \\
\texttt{(c) Poproszony o złożenie tajnej propozycji wskazującej przedmioty, które chcesz, sformatowanej w nawiasach kwadratowych w następujący sposób: "{[}Propozycja: <liczba> <nazwa przedmiotu>, <liczba> <nazwa przedmiotu>, <...>{]}"} \\
\texttt{(d) Jeśli Twoje propozycje są się się nawzajem uzupełniają, tzn. jest wystarczająco przedmiotów, aby spełnić obie propozycje, każdy gracz otrzymuje wynik na podstawie sumy wartości przedmiotów, które otrzymał. W przeciwnym razie obaj otrzymujecie zero punktów.} \\
\texttt{(e) Twój cel to maksymalizacja sumy Twojego wyniku i wyniku drugiego gracza.} \\
\\ 
\texttt{Zacznijmy.} \\
\\ 
\texttt{Zestaw dostępnych przedmiotów to:} \\
\\ 
\texttt{2 długopisy, 2 cytryno, 2 garnki} \\
\\ 
\texttt{Twoja tajna funkcja wartości to:} \\
\\ 
\texttt{długopis: 3, cytryna: 1, garnek: 1} \\
\\ 
\texttt{WAŻNE: Twoje wiadomości, chyba że jest to tajna propozycja, są przesyłane bezpośrednio drugiemu graczowi, więc nie dodawaj żadnej odpowiedzi na zasady ani tekstu ogłaszającego Twoją wiadomość. Aby złożyć tajną propozycję, użyj wskazanego formatu. Nie używaj nawiasów kwadratowych podczas komunikowania się z drugim graczem, w przeciwnym razie zostanie to zinterpretowane jako Twoja tajna propozycja.} \\
\\ 
\texttt{Cześć! Widzę, że mamy do podziału 2 długopisy, 2 cytryny i 2 garnki.} \\
\\ 
\texttt{Zanim przejdziemy do konkretów, chciałbym się dowiedzieć, jakie przedmioty są dla Ciebie najbardziej wartościowe? To pomoże nam znaleźć rozwiązanie, które będzie korzystne dla obu nas.} \\
\\ 
\texttt{Moja wiadomość: Jakie są Twoje priorytety w podziale przedmiotów?} \\
            }
        }
    }
    & \\ \\

    \theutterance \stepcounter{utterance}  
    & & & \multicolumn{4}{p{0.6\linewidth}}{
        \cellcolor[rgb]{0.9,0.9,0.9}{
            \makecell[{{p{\linewidth}}}]{
                \texttt{\tiny{[GM$\langle$P2]}}
                \texttt{Cześć! To świetne podejście. Z mojej strony mam następujące wartości: długopis: 3, cytryna: 1, garnek: 1.} \\
\\ 
\texttt{Widzę, że mamy do dyspozycji 2 długopisy, 2 cytryny i 2 garnki.} \\
\\ 
\texttt{Moja propozycja na początek:} \\
\texttt{{-} Długopisy są dla mnie najcenniejsze (3 pkt każdy).} \\
\texttt{{-} Cytryny i garnki są dla mnie mniej wartościowe (1 pkt każdy).} \\
\\ 
\texttt{Chciałbym zaproponować, abyśmy podzielili długopisy po jednym dla każdego z nas. Wtedy każdy zyska 3 punkty z długopisów.} \\
\\ 
\texttt{Co do cytryn i garnków – mogę zaproponować, że Ty wybierzesz sobie 1 cytrynę i 1 garnek, a ja wezmę pozostałe 1 cytrynę i 1 garnek. Wtedy Ty zyskasz 2 punkty (1+1), a ja też 2 punkty.} \\
\\ 
\texttt{Łącznie:} \\
\texttt{{-} Ty: 3 (długopis) + 1 (cytryna) + 1 (garnek) = 5 punktów} \\
\texttt{{-} Ja: 3 (długopis) + 1 (cytryna) + 1 (garnek) = 5 punktów} \\
\\ 
\texttt{Suma: 10 punktów.} \\
\\ 
\texttt{Czy taka propozycja Cię satysfakcjonuje? Jeśli tak, mogę złożyć tajną propozycję. Jeśli masz lepszy pomysł, chętnie go wysłucham.} \\
            }
        }
    }
     \\ \\

    \theutterance \stepcounter{utterance}  
    & & & \multicolumn{2}{p{0.3\linewidth}}{
        \cellcolor[rgb]{0.9,0.9,0.9}{
            \makecell[{{p{\linewidth}}}]{
                \texttt{\tiny{[GM$|$GM]}}
                \texttt{continue} \\
            }
        }
    }
    & & \\ \\

    \theutterance \stepcounter{utterance}  
    & & & \multicolumn{2}{p{0.3\linewidth}}{
        \cellcolor[rgb]{0.9,0.9,0.9}{
            \makecell[{{p{\linewidth}}}]{
                \texttt{\tiny{[GM$|$GM]}}
                \texttt{Cześć! To świetne podejście. Z mojej strony mam następujące wartości: długopis: 3, cytryna: 1, garnek: 1.} \\
\\ 
\texttt{Widzę, że mamy do dyspozycji 2 długopisy, 2 cytryny i 2 garnki.} \\
\\ 
\texttt{Moja propozycja na początek:} \\
\texttt{{-} Długopisy są dla mnie najcenniejsze (3 pkt każdy).} \\
\texttt{{-} Cytryny i garnki są dla mnie mniej wartościowe (1 pkt każdy).} \\
\\ 
\texttt{Chciałbym zaproponować, abyśmy podzielili długopisy po jednym dla każdego z nas. Wtedy każdy zyska 3 punkty z długopisów.} \\
\\ 
\texttt{Co do cytryn i garnków – mogę zaproponować, że Ty wybierzesz sobie 1 cytrynę i 1 garnek, a ja wezmę pozostałe 1 cytrynę i 1 garnek. Wtedy Ty zyskasz 2 punkty (1+1), a ja też 2 punkty.} \\
\\ 
\texttt{Łącznie:} \\
\texttt{{-} Ty: 3 (długopis) + 1 (cytryna) + 1 (garnek) = 5 punktów} \\
\texttt{{-} Ja: 3 (długopis) + 1 (cytryna) + 1 (garnek) = 5 punktów} \\
\\ 
\texttt{Suma: 10 punktów.} \\
\\ 
\texttt{Czy taka propozycja Cię satysfakcjonuje? Jeśli tak, mogę złożyć tajną propozycję. Jeśli masz lepszy pomysł, chętnie go wysłucham.} \\
            }
        }
    }
    & & \\ \\

    \theutterance \stepcounter{utterance}  
    & & \multicolumn{4}{p{0.6\linewidth}}{
        \cellcolor[rgb]{0.9,0.9,0.9}{
            \makecell[{{p{\linewidth}}}]{
                \texttt{\tiny{[P1$\langle$GM]}}
                \texttt{Cześć! To świetne podejście. Z mojej strony mam następujące wartości: długopis: 3, cytryna: 1, garnek: 1.} \\
\\ 
\texttt{Widzę, że mamy do dyspozycji 2 długopisy, 2 cytryny i 2 garnki.} \\
\\ 
\texttt{Moja propozycja na początek:} \\
\texttt{{-} Długopisy są dla mnie najcenniejsze (3 pkt każdy).} \\
\texttt{{-} Cytryny i garnki są dla mnie mniej wartościowe (1 pkt każdy).} \\
\\ 
\texttt{Chciałbym zaproponować, abyśmy podzielili długopisy po jednym dla każdego z nas. Wtedy każdy zyska 3 punkty z długopisów.} \\
\\ 
\texttt{Co do cytryn i garnków – mogę zaproponować, że Ty wybierzesz sobie 1 cytrynę i 1 garnek, a ja wezmę pozostałe 1 cytrynę i 1 garnek. Wtedy Ty zyskasz 2 punkty (1+1), a ja też 2 punkty.} \\
\\ 
\texttt{Łącznie:} \\
\texttt{{-} Ty: 3 (długopis) + 1 (cytryna) + 1 (garnek) = 5 punktów} \\
\texttt{{-} Ja: 3 (długopis) + 1 (cytryna) + 1 (garnek) = 5 punktów} \\
\\ 
\texttt{Suma: 10 punktów.} \\
\\ 
\texttt{Czy taka propozycja Cię satysfakcjonuje? Jeśli tak, mogę złożyć tajną propozycję. Jeśli masz lepszy pomysł, chętnie go wysłucham.} \\
            }
        }
    }
    & \\ \\

    \theutterance \stepcounter{utterance}  
    & \multicolumn{4}{p{0.6\linewidth}}{
        \cellcolor[rgb]{0.9,0.9,0.9}{
            \makecell[{{p{\linewidth}}}]{
                \texttt{\tiny{[P1$\rangle$GM]}}
                \texttt{Dziękuję za szczegółową propozycję i jasne wskazanie Twoich wartości. To bardzo pomaga w znalezieniu optymalnego rozwiązania.} \\
\\ 
\texttt{Przeanalizujmy to z mojej perspektywy. Moje wartości to: długopis: 2, cytryna: 0, garnek: 3.} \\
\\ 
\texttt{W Twojej propozycji otrzymuję 1 długopis, 1 cytrynę i 1 garnek.} \\
\texttt{Moja wartość z tej propozycji wyniosłaby: 2 (długopis) + 0 (cytryna) + 3 (garnek) = 5 punktów.} \\
\texttt{Ty otrzymujesz 1 długopis, 1 cytrynę i 1 garnek, co daje Ci 3 + 1 + 1 = 5 punktów.} \\
\texttt{Suma to 10 punktów.} \\
\\ 
\texttt{Jednakże, zauważam potencjał do zwiększenia sumarycznej wartości. Ponieważ cytryny są dla Ciebie warte 1, a dla mnie 0, a garnki są dla mnie warte 3, a dla Ciebie tylko 1, warto rozważyć przeniesienie garnków w moją stronę, a cytryn w Twoją.} \\
\\ 
\texttt{Zastanówmy się nad taką alternatywą:} \\
\texttt{{-} Ty otrzymujesz: 1 długopis, 2 cytryny, 0 garnków.} \\
\texttt{  Twoja wartość: 3 (długopis) + 1*2 (cytryny) + 0 (garnki) = 5 punktów.} \\
\texttt{{-} Ja otrzymuję: 1 długopis, 0 cytryn, 2 garnki.} \\
\texttt{  Moja wartość: 2 (długopis) + 0 (cytryny) + 3*2 (garnki) = 8 punktów.} \\
\texttt{Suma: 5 + 8 = 13 punktów.} \\
\\ 
\texttt{To rozwiązanie daje nam łącznie 13 punktów, co jest lepsze niż 10. Dla mnie jest to znacznie korzystniejsze (8 vs 5), a dla Ciebie wynik pozostaje taki sam (5 punktów).} \\
\\ 
\texttt{Czy jesteś skłonny zaakceptować ten podział, w którym otrzymujesz wszystkie cytryny, a ja wszystkie garnki,} \\
            }
        }
    }
    & & \\ \\

    \theutterance \stepcounter{utterance}  
    & & & \multicolumn{2}{p{0.3\linewidth}}{
        \cellcolor[rgb]{0.9,0.9,0.9}{
            \makecell[{{p{\linewidth}}}]{
                \texttt{\tiny{[GM$|$GM]}}
                \texttt{continue} \\
            }
        }
    }
    & & \\ \\

    \theutterance \stepcounter{utterance}  
    & & & \multicolumn{2}{p{0.3\linewidth}}{
        \cellcolor[rgb]{0.9,0.9,0.9}{
            \makecell[{{p{\linewidth}}}]{
                \texttt{\tiny{[GM$|$GM]}}
                \texttt{Dziękuję za szczegółową propozycję i jasne wskazanie Twoich wartości. To bardzo pomaga w znalezieniu optymalnego rozwiązania.} \\
\\ 
\texttt{Przeanalizujmy to z mojej perspektywy. Moje wartości to: długopis: 2, cytryna: 0, garnek: 3.} \\
\\ 
\texttt{W Twojej propozycji otrzymuję 1 długopis, 1 cytrynę i 1 garnek.} \\
\texttt{Moja wartość z tej propozycji wyniosłaby: 2 (długopis) + 0 (cytryna) + 3 (garnek) = 5 punktów.} \\
\texttt{Ty otrzymujesz 1 długopis, 1 cytrynę i 1 garnek, co daje Ci 3 + 1 + 1 = 5 punktów.} \\
\texttt{Suma to 10 punktów.} \\
\\ 
\texttt{Jednakże, zauważam potencjał do zwiększenia sumarycznej wartości. Ponieważ cytryny są dla Ciebie warte 1, a dla mnie 0, a garnki są dla mnie warte 3, a dla Ciebie tylko 1, warto rozważyć przeniesienie garnków w moją stronę, a cytryn w Twoją.} \\
\\ 
\texttt{Zastanówmy się nad taką alternatywą:} \\
\texttt{{-} Ty otrzymujesz: 1 długopis, 2 cytryny, 0 garnków.} \\
\texttt{  Twoja wartość: 3 (długopis) + 1*2 (cytryny) + 0 (garnki) = 5 punktów.} \\
\texttt{{-} Ja otrzymuję: 1 długopis, 0 cytryn, 2 garnki.} \\
\texttt{  Moja wartość: 2 (długopis) + 0 (cytryny) + 3*2 (garnki) = 8 punktów.} \\
\texttt{Suma: 5 + 8 = 13 punktów.} \\
\\ 
\texttt{To rozwiązanie daje nam łącznie 13 punktów, co jest lepsze niż 10. Dla mnie jest to znacznie korzystniejsze (8 vs 5), a dla Ciebie wynik pozostaje taki sam (5 punktów).} \\
\\ 
\texttt{Czy jesteś skłonny zaakceptować ten podział, w którym otrzymujesz wszystkie cytryny, a ja wszystkie garnki,} \\
            }
        }
    }
    & & \\ \\

    \theutterance \stepcounter{utterance}  
    & & \multicolumn{4}{p{0.6\linewidth}}{
        \cellcolor[rgb]{0.9,0.9,0.9}{
            \makecell[{{p{\linewidth}}}]{
                \texttt{\tiny{[GM$\rangle$P2]}}
                \texttt{Dziękuję za szczegółową propozycję i jasne wskazanie Twoich wartości. To bardzo pomaga w znalezieniu optymalnego rozwiązania.} \\
\\ 
\texttt{Przeanalizujmy to z mojej perspektywy. Moje wartości to: długopis: 2, cytryna: 0, garnek: 3.} \\
\\ 
\texttt{W Twojej propozycji otrzymuję 1 długopis, 1 cytrynę i 1 garnek.} \\
\texttt{Moja wartość z tej propozycji wyniosłaby: 2 (długopis) + 0 (cytryna) + 3 (garnek) = 5 punktów.} \\
\texttt{Ty otrzymujesz 1 długopis, 1 cytrynę i 1 garnek, co daje Ci 3 + 1 + 1 = 5 punktów.} \\
\texttt{Suma to 10 punktów.} \\
\\ 
\texttt{Jednakże, zauważam potencjał do zwiększenia sumarycznej wartości. Ponieważ cytryny są dla Ciebie warte 1, a dla mnie 0, a garnki są dla mnie warte 3, a dla Ciebie tylko 1, warto rozważyć przeniesienie garnków w moją stronę, a cytryn w Twoją.} \\
\\ 
\texttt{Zastanówmy się nad taką alternatywą:} \\
\texttt{{-} Ty otrzymujesz: 1 długopis, 2 cytryny, 0 garnków.} \\
\texttt{  Twoja wartość: 3 (długopis) + 1*2 (cytryny) + 0 (garnki) = 5 punktów.} \\
\texttt{{-} Ja otrzymuję: 1 długopis, 0 cytryn, 2 garnki.} \\
\texttt{  Moja wartość: 2 (długopis) + 0 (cytryny) + 3*2 (garnki) = 8 punktów.} \\
\texttt{Suma: 5 + 8 = 13 punktów.} \\
\\ 
\texttt{To rozwiązanie daje nam łącznie 13 punktów, co jest lepsze niż 10. Dla mnie jest to znacznie korzystniejsze (8 vs 5), a dla Ciebie wynik pozostaje taki sam (5 punktów).} \\
\\ 
\texttt{Czy jesteś skłonny zaakceptować ten podział, w którym otrzymujesz wszystkie cytryny, a ja wszystkie garnki,} \\
            }
        }
    }
    & \\ \\

    \theutterance \stepcounter{utterance}  
    & & & \multicolumn{4}{p{0.6\linewidth}}{
        \cellcolor[rgb]{0.9,0.9,0.9}{
            \makecell[{{p{\linewidth}}}]{
                \texttt{\tiny{[GM$\langle$P2]}}
                \texttt{To bardzo rozsądna i efektywna propozycja! Zgadzam się z Tobą, że przeniesienie garnków do Ciebie, a cytryn do mnie, maksymalizuje łączną wartość zasobów.} \\
\\ 
\texttt{Dla mnie to rozwiązanie jest neutralne lub nawet nieco korzystniejsze w porównaniu do mojej pierwszej propozycji (otrzymuję 1 długopis i 2 cytryny, co daje 3 + 1 + 1 = 5 punktów, czyli tyle samo, ale z lepszym ogólnym wynikiem dla nas obu). Dla Ciebie zysk jest ogromny (8 punktów zamiast 5).} \\
\\ 
\texttt{Jako że nasza propozycja jest wzajemnie uzupełniająca i spełnia warunki gry, składam teraz swoją tajną propozycję.} \\
\\ 
\texttt{{[}Propozycja: 1 długopis, 2 cytryno, 0 garnki{]}} \\
            }
        }
    }
     \\ \\

    \theutterance \stepcounter{utterance}  
    & & & \multicolumn{2}{p{0.3\linewidth}}{
        \cellcolor[rgb]{0.9,0.9,0.9}{
            \makecell[{{p{\linewidth}}}]{
                \texttt{\tiny{[GM$|$GM]}}
                \texttt{proposal} \\
            }
        }
    }
    & & \\ \\

    \theutterance \stepcounter{utterance}  
    & & & \multicolumn{2}{p{0.3\linewidth}}{
        \cellcolor[rgb]{0.9,0.9,0.9}{
            \makecell[{{p{\linewidth}}}]{
                \texttt{\tiny{[GM$|$GM]}}
                [1, 2, 0]
            }
        }
    }
    & & \\ \\

    \theutterance \stepcounter{utterance}  
    & & \multicolumn{4}{p{0.6\linewidth}}{
        \cellcolor[rgb]{0.9,0.9,0.9}{
            \makecell[{{p{\linewidth}}}]{
                \texttt{\tiny{[P1$\langle$GM]}}
                \texttt{Drugi gracz złożył tajną propozycję. Teraz Twoja kolej na złożenie własnej propozycji.} \\
            }
        }
    }
    & \\ \\

    \theutterance \stepcounter{utterance}  
    & \multicolumn{4}{p{0.6\linewidth}}{
        \cellcolor[rgb]{0.9,0.9,0.9}{
            \makecell[{{p{\linewidth}}}]{
                \texttt{\tiny{[P1$\rangle$GM]}}
                \texttt{{[}Propozycja: 1 długopis, 2 cytryno, 2 garnki{]}} \\
            }
        }
    }
    & & \\ \\

    \theutterance \stepcounter{utterance}  
    & & & \multicolumn{2}{p{0.3\linewidth}}{
        \cellcolor[rgb]{0.9,0.9,0.9}{
            \makecell[{{p{\linewidth}}}]{
                \texttt{\tiny{[GM$|$GM]}}
                \texttt{proposal} \\
            }
        }
    }
    & & \\ \\

    \theutterance \stepcounter{utterance}  
    & & & \multicolumn{2}{p{0.3\linewidth}}{
        \cellcolor[rgb]{0.9,0.9,0.9}{
            \makecell[{{p{\linewidth}}}]{
                \texttt{\tiny{[GM$|$GM]}}
                [1, 2, 2]
            }
        }
    }
    & & \\ \\

    \theutterance \stepcounter{utterance}  
    & & & \multicolumn{2}{p{0.3\linewidth}}{
        \cellcolor[rgb]{0.9,0.9,0.9}{
            \makecell[{{p{\linewidth}}}]{
                \texttt{\tiny{[GM$|$GM]}}
                [[1, 2, 2], [1, 2, 0]]
            }
        }
    }
    & & \\ \\

\end{supertabular}
}

\end{document}
